
A self-supervised model trained on Waterbirds achieves 62.10\% average accuracy --- yet its worst-group accuracy collapses to \textbf{8.57\%}. The same images, fed to a supervised model, yield \textbf{81.93\%} worst-group accuracy. This 73.36 pp gap is not a failure of representation capacity. It is a failure of geometry.

Self-supervised learning (SSL) has become the dominant paradigm for pretraining visual representations at scale. SimCLR, MoCo, and DINO produce features that transfer broadly across downstream tasks --- yet on spurious correlation benchmarks, these same models fail catastrophically on minority groups while appearing healthy by average accuracy. A DINO model that correctly identifies 97.25\% of waterbirds on water backgrounds identifies only 8.57\% of waterbirds on land backgrounds. The majority-group accuracy provides a misleading certificate of success while hiding a systematic failure that real-world deployment would expose.

The standard view frames this as a representational problem: SSL encodes spurious correlations (background texture) rather than core features (bird morphology) [Zhang \& Ré, 2022; Yadav et al., 2026]. From this view, the fix is post-hoc: resample or reweight training examples based on loss-based proxies \citep{ghaznavi2023lfr}, or diversify augmentations to reduce the informativeness of spurious features \citep{yadav2026crossvariant}. These strategies operate on the representation as given, not on the training process that created it.

We propose a different lens: the SSL shortcut problem is geometric. When SGD minimizes the InfoNCE objective on spurious correlation benchmarks, the loss landscape develops \textit{directional curvature asymmetry} --- sharper curvature along spurious feature directions than along core feature directions. This sharpness anisotropy is the geometric signature by which spurious features become dominant in the learned representation. Post-hoc interventions that operate on fixed representations cannot undo this geometric structure. An intervention must happen during training, at the optimizer level.

\textbf{Our key insight:} SSL-SGD creates an anisotropic loss landscape where spurious feature directions are geometrically privileged --- sharper, more sensitive to perturbation, and therefore more reliably encoded. Sharpness-Aware Minimization (SAM) \citep{foret2021sam}, by seeking flat minima, preferentially flattens these spurious-direction curvature peaks. Crucially, the spurious directions can be identified \textit{without group labels} using a linear probe loss variance proxy \citep{ghaznavi2023lfr}, making the entire pipeline annotation-free.

This geometric reframing unifies two previously disconnected lines of work: the SAM literature (which analyzes flat minima but has not, to our knowledge, been applied to SSL spurious correlations) and the spurious correlation literature (which has not used geometric loss landscape analysis as a diagnostic). Building on Gatmiry et al.'s [2024] theoretical link between SAM and simplicity bias in supervised settings, and G2-SAM's \citep{park2025g2sam} demonstration that group-wise sharpness reduction improves worst-group accuracy, we ask: does a directly analogous geometric mechanism operate in SSL with InfoNCE objectives?

We make the following contributions:

\textbf{1. Empirical confirmation of the SSL shortcut severity.} We document an 8.57\% WGA for DINO SSL versus 81.93\% WGA for supervised training (73.36 pp gap) on Waterbirds, using identical evaluation protocols. This gap persists throughout training --- SSL average accuracy converges normally while WGA stagnates below 25\% --- confirming that the shortcut encoding is geometrically stable, not a transient artifact of early training.

\textbf{2. Geometric hypothesis and measurement protocol (designed, pending validation).} We introduce an explicit geometric hypothesis for SSL shortcut encoding: that Hessian sharpness anisotropy along spurious vs. random feature directions exceeds 1.$2\times$ in converged SSL models, and correlates negatively with worst-group accuracy (Pearson r < $-0.5)$ across training checkpoints. We design an annotation-free measurement protocol combining SAM perturbations and the LFR loss-variance proxy.

\textbf{3. SAM as a training-time geometric intervention (designed, pending validation).} We propose applying SAM during SSL pretraining to flatten the spurious-direction curvature peaks identified by the geometric diagnostic. Unlike post-hoc methods, SAM intervenes during the geometric formation of the representation, addressing the cause rather than the symptom.

Contributions 2 and 3 are complete research protocols with pending experimental validation from 200-epoch runs; they are presented alongside the confirmed finding to enable community engagement with the geometric hypothesis and the measurement methodology.

The remainder of the paper is organized as follows. Section 2 reviews related work on group robustness, annotation-free debiasing, and geometric loss landscape analysis. Section 3 details our methodology, including the anisotropy measurement protocol and the SAM-SSL training procedure. Section 4 describes our experimental setup. Section 5 presents results, including the confirmed WGA gap and the status of ongoing anisotropy measurement. Section 6 discusses implications and limitations. Section 7 concludes.

